\documentclass{book}
\usepackage[utf8]{inputenc}
\usepackage[T1]{fontenc}
\usepackage{lmodern}        % Fuentes vectoriales Type 1
\usepackage{microtype}      % Mejora la legibilidad
\usepackage[spanish]{babel}
\usepackage[margin=2.5cm]{geometry}
\usepackage{amsmath}
\usepackage{amssymb}
\usepackage{graphicx}
\usepackage{booktabs}
\usepackage{tabularx}
\usepackage{sidecap}
\usepackage{url}

\graphicspath{{imagenes/}}

\AtBeginDocument{\shorthandoff{"}}

\begin{document}

% ==========================================
% PÁGINA DE DETALLES DE LA SERIE Y COMITÉ EDITORIAL
% ==========================================
\thispagestyle{empty}
\noindent\textbf{FÍSICA BIOLÃ``GICA Y MÉDICA, INGENIERÍA BIOMÉDICA}

\vspace{1em}
\noindent Los campos de la física biológica y médica y de la ingeniería biomédica son amplios, multidisciplinarios y dinámicos. Se encuentran en la encrucijada de la investigación de frontera en física, biología, química y medicina. La serie \textit{Biological and Medical Physics, Biomedical Engineering} está diseñada para ser integral, cubriendo una amplia gama de temas importantes para el estudio de las ciencias físicas, químicas y biológicas. Su objetivo es proporcionar a científicos e ingenieros libros de texto, monografías y obras de referencia para abordar la creciente necesidad de información.

\noindent Los libros de la serie enfatizan áreas establecidas y emergentes de la ciencia, incluyendo la biofísica molecular, de membranas y matemática; la captación y conversión de energía fotosintética; el procesamiento de información; los principios físicos de la genética; las comunicaciones sensoriales; las redes de autómatas, redes neuronales y autómatas celulares. De igual importancia será la cobertura de aspectos aplicados de la física biológica y médica y de la ingeniería biomédica, tales como componentes y dispositivos electrónicos moleculares, biosensores, medicina, imágenes, principios físicos de la producción de energía renovable, prótesis avanzadas, y control e ingeniería ambiental.

\vspace{1.5em}
\noindent\textbf{Editor en Jefe:} \\
Elias Greenbaum, Oak Ridge National Laboratory, Oak Ridge, Tennessee, EE. UU.

\vspace{1.5em}
\noindent\textbf{Comité Editorial:}
\begin{small}
\begin{itemize}
    \item \textbf{Masuo Aizawa}, Departamento de Bioingeniería, Instituto de Tecnología de Tokio, Yokohama, Japón
    \item \textbf{Olaf S. Andersen}, Departamento de Fisiología, Biofísica y Medicina Molecular, Universidad de Cornell, Nueva York, EE. UU.
    \item \textbf{Robert H. Austin}, Departamento de Física, Universidad de Princeton, Princeton, Nueva Jersey, EE. UU.
    \item \textbf{James Barber}, Departamento de Bioquímica, Imperial College of Science, Technology and Medicine, Londres, Inglaterra
    \item \textbf{Howard C. Berg}, Departamento de Biología Molecular y Celular, Universidad de Harvard, Cambridge, Massachusetts, EE. UU.
    \item \textbf{Victor Bloomfield}, Departamento de Bioquímica, Universidad de Minnesota, St. Paul, Minnesota, EE. UU.
    \item \textbf{Robert Callender}, Departamento de Bioquímica, Albert Einstein College of Medicine, Bronx, Nueva York, EE. UU.
    \item \textbf{Britton Chance}, Departamento de Bioquímica/Biofísica, Universidad de Pensilvania, Filadelfia, EE. UU.
    \item \textbf{Steven Chu}, Lawrence Berkeley National Laboratory, Berkeley, California, EE. UU.
    \item \textbf{Louis J. DeFelice}, Departamento de Farmacología, Universidad de Vanderbilt, Nashville, Tennessee, EE. UU.
    \item \textbf{Johann Deisenhofer}, Instituto Médico Howard Hughes, Universidad de Texas, Dallas, Texas, EE. UU.
    \item \textbf{George Feher}, Departamento de Física, Universidad de California, San Diego, La Jolla, California, EE. UU.
    \item \textbf{Hans Frauenfelder}, Los Alamos National Laboratory, Los Alamos, Nuevo México, EE. UU.
    \item \textbf{Ivar Giaever}, Rensselaer Polytechnic Institute, Troy, Nueva York, EE. UU.
    \item \textbf{Sol M. Gruner}, Universidad de Cornell, Ithaca, Nueva York, EE. UU.
    \item \textbf{Judith Herzfeld}, Departamento de Química, Universidad de Brandeis, Waltham, Massachusetts, EE. UU.
    \item \textbf{Mark S. Humayun}, Doheny Eye Institute, Los Ángeles, California, EE. UU.
    \item \textbf{Pierre Joliot}, Institut de Biologie Physico-Chimique, Fondation Edmond de Rothschild, París, Francia
    \item \textbf{Lajos Keszthelyi}, Instituto de Biofísica, Academia de Ciencias de Hungría, Szeged, Hungría
    \item \textbf{Robert S. Knox}, Departamento de Física y Astronomía, Universidad de Rochester, Rochester, Nueva York, EE. UU.
    \item \textbf{Aaron Lewis}, Departamento de Física Aplicada, Universidad Hebrea, Jerusalén, Israel
    \item \textbf{Stuart M. Lindsay}, Departamento de Física y Astronomía, Universidad Estatal de Arizona, Tempe, Arizona, EE. UU.
    \item \textbf{David Mauzerall}, Universidad Rockefeller, Nueva York, Nueva York, EE. UU.
    \item \textbf{Eugenie V. Mielczarek}, Departamento de Física y Astronomía, Universidad George Mason, Fairfax, Virginia, EE. UU.
    \item \textbf{Markolf Niemz}, Facultad de Medicina de Mannheim, Universidad de Heidelberg, Mannheim, Alemania
    \item \textbf{V. Adrian Parsegian}, Physical Science Laboratory, National Institutes of Health, Bethesda, Maryland, EE. UU.
    \item \textbf{Linda S. Powers}, Universidad de Arizona, Tucson, Arizona, EE. UU.
    \item \textbf{Earl W. Prohofsky}, Departamento de Física, Universidad de Purdue, West Lafayette, Indiana, EE. UU.
    \item \textbf{Andrew Rubin}, Departamento de Biofísica, Universidad Estatal de Moscú, Moscú, Rusia
    \item \textbf{Michael Seibert}, National Renewable Energy Laboratory, Golden, Colorado, EE. UU.
    \item \textbf{David Thomas}, Departamento de Bioquímica, Facultad de Medicina de la Universidad de Minnesota, Minneapolis, Minnesota, EE. UU.
\end{itemize}
\end{small}

\vspace{1.5em}
\noindent Más información sobre esta serie en \url{http://www.springer.com/series/3740}
\clearpage

% ==========================================
% PORTADA INTERNA (TITLE PAGE)
% ==========================================
\thispagestyle{empty}
\begin{center}
    \vspace*{3cm}
    {\large\textbf{Física Biológica y Médica, Ingeniería Biomédica}}
    
    \vspace{2.5cm}
    {\Huge\textbf{Irving P. Herman}}
    
    \vspace{1.5cm}
    {\Huge\textbf{Física del Cuerpo Humano}}
    
    \vspace{1.2cm}
    {\Large\textbf{Segunda Edición}}
    
    \vspace*{\fill}
    {\large\textbf{Springer}}
\end{center}
\clearpage

% ==========================================
% PÁGINA DE COPYRIGHT
% ==========================================
\thispagestyle{empty}
\noindent\textbf{Irving P. Herman} \\
Departamento de Física Aplicada \\
y Matemáticas Aplicadas \\
Universidad de Columbia \\
Nueva York, NY, EE. UU.

\vspace{3em}
\noindent ISSN 1618-7210 \hfill ISSN 2197-5647 (electrónico) \\
Biological and Medical Physics, Biomedical Engineering \\
ISBN 978-3-319-23930-9 \hfill ISBN 978-3-319-23932-3 (eBook) \\
DOI 10.1007/978-3-319-23932-3

\vspace{1.5em}
\noindent Número de Control de la Biblioteca del Congreso: 2015951417

\vspace{1.5em}
\noindent Springer Cham Heidelberg New York Dordrecht London

\vspace{1.5em}
\noindent \copyright\ Springer International Publishing Switzerland 2006, 2016

\vspace{1.5em}
\noindent Esta obra está sujeta a derechos de autor. Todos los derechos están reservados por la Editorial, ya sea que se trate de la totalidad o parte del material, específicamente los derechos de traducción, reimpresión, reutilización de ilustraciones, recitación, difusión, reproducción en microfilmes o en cualquier otra forma física, y transmisión o almacenamiento de información y recuperación, adaptación informática, software de computadora, o por metodología similar o diferente conocida ahora o desarrollada en el futuro.

\vspace{1em}
\noindent El uso de nombres descriptivos generales, nombres registrados, marcas comerciales, marcas de servicio, etc. en esta publicación no implica, incluso en ausencia de una declaración específica, que dichos nombres estén exentos de las leyes y reglamentos de protección pertinentes y, por lo tanto, libres para su uso general.

\vspace{1em}
\noindent La editorial, los autores y los editores pueden asumir con seguridad que los consejos e información de este libro se consideran verdaderos y precisos a la fecha de publicación. Ni la editorial ni los autores o editores ofrecen una garantía, expresa o implícita, con respecto al material aquí contenido o por cualquier error u omisión que pudiera haberse cometido.

\vspace{2.5em}
\noindent Impreso en papel libre de ácido.

\vspace{1.5em}
\noindent Springer International Publishing AG Switzerland es parte de Springer Science+Business Media (www.springer.com)
\clearpage

% ==========================================
% DEDICATORIA
% ==========================================
\thispagestyle{empty}
\vspace*{7cm}
\begin{center}
    \itshape
    Este libro está dedicado a Daniel, Jonathan y Janet, \\
    a Sandy Koufax, \\
    y a la vida de Ron McNair.
\end{center}
\clearpage

% ==========================================
% PREFACIO A LA SEGUNDA EDICIÃ``N
% ==========================================
\chapter*{Prefacio a la Segunda Edición}
\addcontentsline{toc}{chapter}{Prefacio a la Segunda Edición}

Cerca del final de uno de los discursos más famosos de la historia estadounidense, Abraham Lincoln hizo referencia al «gobierno del pueblo, por el pueblo, para el pueblo» \cite{2nd-5}. En un sentido similar, un texto sobre la física y el cuerpo humano podría cubrir la «física del cuerpo, por el cuerpo y para el cuerpo», pero tal vez eso sería demasiado amplio. Nuestro enfoque principal en esta segunda edición sigue siendo la física del cuerpo, concretamente la física de la fisiología, y el texto está nuevamente organizado en gran medida por temas de fisiología. Sigue habiendo un fuerte énfasis en la modelación de la fisiología. Como se señala en el excelente libro de Rideout \cite{2nd-7}, la modelación general en fisiología incluye la del momento (como el flujo sanguíneo), la masa (oxígeno, dióxido de carbono, etc., en la sangre), la energía (flujo de calor) y la información (señales en los nervios, hormonas en el torrente sanguíneo) \cite{2nd-7}. Los modelos que acoplan múltiples sistemas también son importantes. (Después de todo, todos somos múltiples sistemas fisiológicos fuertemente acoplados que, por supuesto, están determinados en última instancia por la física). Seguiremos limitándonos a los aspectos más físicos y menos químicos de la modelación fisiológica; la modelación química se aborda adecuadamente en otras fuentes \cite{2nd-1, 2nd-2, 2nd-3, 2nd-4, 2nd-6, 2nd-7}. Dentro de nuestro enfoque de «física del cuerpo», hemos mejorado la discusión sobre la física de las enfermedades, el envejecimiento, el movimiento y la estructura y función de los músculos, y ahora exploramos cómo el equilibrio, el metabolismo y el sistema circulatorio cambian durante el embarazo y el parto. La física de los eventos transversales del ciclo de vida, como el embarazo, las lesiones, las enfermedades y trastornos, el envejecimiento y la muerte, trasciende los temas de cualquier capítulo individual. Ahora se resumen en el nuevo Apéndice F; este apéndice también describe cómo localizar las discusiones sobre estos temas dentro del texto principal.

La «física por el cuerpo» incluye la física de la locomoción común y la que ocurre durante eventos atléticos y colisiones. Estos temas se abordaron en varias partes de la primera edición, incluyendo muchos problemas. Ahora se discuten con más detalle tanto en el texto como en los problemas. En particular, la física del movimiento, la kinesiología y el equilibrio \cite{2nd-8} se discuten de manera más directa mediante el análisis del modelado multisegmento y el tratamiento ampliado de la física del deporte. La discusión sobre el rendimiento humano en tareas cotidianas y en los extremos de la competencia atlética también trasciende los temas de cualquier capítulo individual y acopla muchos de ellos, especialmente los que abordan la estática, el movimiento, el metabolismo, la respiración y la circulación sanguínea. Para ayudar a unificar la discusión sobre el atletismo, se ha añadido el Apéndice G para vincular los temas relevantes en estos capítulos y unificar las referencias a diferentes actividades deportivas.

La «física para el cuerpo» incluye la física del análisis de mediciones y el diagnóstico, así como la física del tratamiento, gran parte de la cual se denomina física médica. La primera edición contenía cierta discusión sobre mediciones y diagnóstico, y prácticamente nada sobre el tratamiento. La presentación del análisis de mediciones incluyó varios métodos y un poco sobre los resultados de las mediciones. En el nuevo Apéndice H se presenta ahora una descripción general de este tema, que incluye un localizador que describe dónde se detalla el análisis físico del cuerpo en el texto y en los problemas. Aparte de la nueva y breve discusión en el Apéndice H, el tratamiento sigue sin abordarse en esta edición.

Este texto aún puede adaptarse para su uso en cursos de pregrado de nivel inferior o superior, así como para algunos cursos de posgrado de nivel inicial. Puede cumplir las funciones de texto principal para un curso completo o servir como fuente secundaria, por ejemplo, mediante el uso de capítulos individuales. Los Apéndices F, G y H unifican los temas transversales presentados en varios capítulos, y pueden utilizarse para guiar la discusión de estos temas.

Esta edición incluye numerosos problemas adicionales que abordan el material del texto y lo exploran con mucho más detalle. Ahora se proporcionan muchas soluciones adicionales al final del libro. El autor recomienda encarecidamente que, al leer cualquier sección del texto, al menos se examine el contenido de los problemas relacionados al final del capítulo (incluso si no se planea resolverlos), ya que muchos enfoques e ideas esenciales se tratan únicamente en los problemas.

En la edición anterior se citaron varios textos excelentes sobre temas relacionados. En particular, destacamos nuevamente varios de los libros sobresalientes que describen modelos matemáticos de fisiología \cite{2nd-1, 2nd-2, 2nd-3, 2nd-4, 2nd-6, 2nd-7}, los cuales están en su mayor parte a un nivel más avanzado que el de este texto. Además, resulta de interés más general la obra reciente «The Story of the Human Body» de Daniel E. Lieberman, que analiza la evolución del cuerpo humano e incluye perspectivas extraordinarias sobre la estructura y función de los cuerpos de los humanos modernos.

El autor está en deuda con Julio Fernandez, Glenn Fleisig, Joel Gersten, Daniel Herman, Jonathan Herman, Ronald Koder, Jeffrey Kysar, David Lerner, Kristin Myers, Herman Pontzer y Alejandro Tello, y muchos otros por sus comentarios, sugerencias y aportes mientras preparaba esta segunda edición; con los numerosos estudiantes, incluyendo a Sarah Beck, que encontraron errores en la primera edición; y con quienes ayudaron en la preparación de la primera edición. También le gustaría agradecer al sistema de bibliotecas de la Universidad de Columbia por la rápida recuperación de libros de referencia y artículos (y por permitirle al autor renovar los libros muchas, muchas veces).

\vspace{1.5em}
\noindent Nueva York, NY \hfill Irving P. Herman \\
Julio de 2015

\vspace{2em}
\begin{thebibliography}{9}
\bibitem{2nd-1} E. Carson, C. Cobelli, \textit{Modelling Methodology for Physiology and Medicine} (Academic, San Diego, 2001)
\bibitem{2nd-2} J.L. Hargrove, \textit{Dynamic Modeling in the Health Sciences} (Springer, Berlín, 1998)
\bibitem{2nd-3} F.C. Hoppensteadt, C.S. Peskin, \textit{Modeling and Simulation in Medicine and the Life Sciences}, 2.ª ed. (Springer, Berlín, 2002) (La primera edición de 1992 se llamó \textit{Mathematics in Medicine and the Life Sciences})
\bibitem{2nd-4} J. Keener, J. Sneyd (eds.), \textit{Mathematical Physiology II: Systems Physiology}, 2.ª ed. (Springer, Berlín, 2009)
\bibitem{2nd-5} A. Lincoln, \textit{Lincoln: Speeches and Writings: 1859--1865} (Library of America, Nueva York, 2012)
\bibitem{2nd-6} J.T. Ottesen, M.S. Olufsen, J.K. Larsen (eds.), \textit{Applied Mathematical Models in Human Physiology} (SIAM, Filadelfia, 2004)
\bibitem{2nd-7} V.C. Rideout, \textit{Mathematical and Computer Modeling of Physiological Systems} (Prentice-Hall, Englewood Cliffs, 1991)
\bibitem{2nd-8} N. Hamilton, W. Weimar, K. Luttgens, \textit{Kinesiology: Scientific Basis of Human Motion}, 12.ª ed. (McGraw-Hill, Nueva York, 2012)
\end{thebibliography}

% Control de Calidad: 7 párrafos origen -> 7 párrafos salida. Mapeo 1:1 verificado. Sin omisiones.

\clearpage

% ==========================================
% PREFACIO A LA PRIMERA EDICIÃ``N
% ==========================================
\chapter*{Prefacio a la Primera Edición}
\addcontentsline{toc}{chapter}{Prefacio a la Primera Edición}

La física explica todo desde el principio hasta el final de cualquier descripción completa del cuerpo humano. Una discusión tan exhaustiva debería comenzar con la estructura básica de la materia, tal como la explica la mecánica cuántica ---la física a pequeñas dimensiones---, y terminar con la mecánica del movimiento humano, la energética del metabolismo, la dinámica de fluidos del flujo sanguíneo a través de los vasos, los mecanismos para hablar y oír, y el sistema de imágenes ópticas que llamamos el ojo. Todas las combinaciones requeridas de átomos para formar las complejas moléculas y órganos de los organismos que viven y se reproducen pueden explicarse mediante la mecánica cuántica; sin embargo, tales explicaciones pueden volverse bastante complejas. Los campos de la química y la biología se han desarrollado, en parte, para explicar la brecha entre los extremos: la microfísica y la macrofísica de organismos como el cuerpo humano.

Este libro se enfoca principalmente en el extremo de la macrofísica del cuerpo humano. Asumiremos que los átomos forman moléculas que forman células que forman órganos. Queremos comprender la física de los órganos humanos y de los propios seres humanos. Aplicaremos y ampliaremos en cierta medida la física de primer año de universidad para ver cómo funciona el cuerpo. Además de aplicar conceptos físicos al cuerpo, intentaremos comprender el cuerpo desde un punto de vista que es más numérico del que se suele adoptar en las presentaciones biológicas y médicas.

Una forma de caracterizar este texto es diciendo lo que es y lo que no es. Ciertamente trata sobre la física del cuerpo humano. No trata sobre anatomía humana, aunque necesitaremos utilizar algunos conceptos anatómicos básicos. No trata sobre fisiología humana, aunque puede llamarse un libro sobre la física de la fisiología. No es una monografía en ingeniería biomédica en sí, aunque aproximadamente la mitad de este volumen se refiere a la biomecánica, un área importante en la ingeniería biomédica. La física médica está más estrechamente relacionada con la física de la salud, el uso de radiación ionizante, las imágenes y la instrumentación, que con la macrofísica del cuerpo. La biofísica se refiere a cómo se puede utilizar la física para estudiar la biología y se centra mucho más en la base molecular y celular de lo que lo haremos nosotros (ver el Apéndice E). Se podría decir que la física del cuerpo humano es sinónimo de comprender la máquina humana.

Nuestro objetivo es comprender los aspectos físicos relacionados con el cuerpo humano, en parte mediante la resolución de problemas para profundizar en esta comprensión. El enfoque no consiste en absoluto en aprender y memorizar terminología médica. Aún así, se introducirán y utilizarán algunos conceptos muy básicos de anatomía y fisiología. Varios de los muchos y excelentes textos generales de anatomía y fisiología se citan al final \cite{1st-11, 1st-16, 1st-21, 1st-22, 1st-23, 1st-24, 1st-25, 1st-26, 1st-27, 1st-29}.

Un tema recurrente a lo largo de este texto es el desarrollo y posterior uso de modelos simples y, posteriormente, más refinados de la macrofísica del cuerpo humano \cite{1st-7, 1st-13, 1st-15}. Los físicos tienden a modelar los conceptos en los términos más simples posibles al principio. Por ejemplo, a orden cero, un físico modelaría una vaca como una esfera. (Esto se utiliza a veces como parte de un chiste). Aquí nos volveremos un poco más complejos, pero no mucho más. Otro tema consiste en abordar de manera cuantitativa cuestiones de la biología humana que a menudo se tratan solo cualitativamente. El llamado de Burton a un pensamiento más cuantitativo en fisiología en \textit{Physiology by Numbers} \cite{1st-5} es muy apreciado por el autor. Además, presentaremos datos fisiológicos reales y los vincularemos con análisis y modelado cuantitativos.

Si hay una fuerza aplicada, energía, flujo de fluidos, un rayo de luz, una corriente eléctrica, o un campo eléctrico o magnético asociado con el cuerpo, lo llamaremos física y lo analizaremos. Tenderemos a evitar temas que profundicen más en cuestiones de química y biología, pero abordaremos brevemente cuestiones de fisicoquímica que involucran gradientes de concentración y similares, en la medida en que se relacionen con el intercambio de fluidos en los capilares y la conducción en los nervios. Aunque enfatizamos la física del cuerpo por encima de la instrumentación utilizada para realizar mediciones físicas en el cuerpo y explorar la función corporal, dicha instrumentación se aborda según sea necesario.

Nuestra intención es utilizar la física básica y no enseñarla, en particular desde cero. Muchos capítulos incluyen una breve revisión de los principios físicos necesarios en ese capítulo y en los capítulos posteriores. Algunos temas se desarrollan un poco más, y algunos incluso un poco más allá, y estos se identifican como temas avanzados. Se ofrecen descripciones más detalladas para temas que rara vez se cubren en detalle durante un curso de física de dos semestres, como los fluidos (Cap. 7), la acústica (Cap. 10) y la óptica (Cap. 11), y para áreas utilizadas en varios contextos, como el movimiento armónico (Cap. 3). Se utiliza algo de cálculo diferencial e integral. (La diferenciación parcial se usa con moderación, y principalmente en las secciones etiquetadas como temas avanzados). En el Apéndice C se presenta una breve revisión de las soluciones a las ecuaciones diferenciales simples utilizadas aquí para ayudar a los estudiantes con una formación limitada en cálculo.

Comenzaremos con una comparación de la terminología médica y física en el Cap. 1. El primer capítulo también incluye una discusión sobre el ser humano «estándar» e introduce el concepto de relaciones de escala. Podemos agrupar los temas de los siguientes 11 capítulos en cuatro áreas de la física del cuerpo humano: (1) En los Caps. 2 al 5, se analiza la mecánica del cuerpo estático (Cap. 2) y del cuerpo en movimiento (Cap. 3), que luego se vinculan con las propiedades mecánicas de los materiales del cuerpo (Cap. 4) y los motores del cuerpo: los músculos (Cap. 5); estos temas pueden caracterizarse como \textit{Locomoción en tierra}. (2) El segundo tema, \textit{Energética del metabolismo corporal}, se analiza en el Cap. 6 y es necesario para comprender las discusiones sobre la locomoción y la función del cuerpo que lo preceden y lo suceden. (3) Los Capítulos 7 al 9 cubren la \textit{Locomoción de los humanos en fluidos} (distinta de la locomoción en tierra) y el \textit{Movimiento de fluidos en los humanos}. El Capítulo 7 ofrece una visión general de la física de los fluidos y aborda la locomoción en el agua (natación) y en el aire (sobre el suelo, al menos la perspectiva del vuelo humano). Los Capítulos 8 y 9 cubren respectivamente la fluídica de la sangre (sistema cardiovascular) y del aire (sistema respiratorio) en el cuerpo. (4) El Capítulo 10 explora la acústica de las ondas sonoras en el habla y la audición. La óptica de los ojos y la visión se investiga en el Cap. 11. Las propiedades eléctricas básicas del cuerpo se desarrollan en el Cap. 12, junto con una breve descripción de las propiedades magnéticas del cuerpo. Por lo tanto, estos tres capítulos abordan respectivamente las ondas sonoras, electromagnéticas y eléctricas, que colectivamente podemos llamar \textit{Ondas y señales}. (La naturaleza electromagnética de las ondas de luz no se analiza en el Cap. 11).

El Capítulo 13 examina cómo el cuerpo utiliza automáticamente el principio de ingeniería básico de retroalimentación y control en la regulación de todos los aspectos de su función. Se describe la física de la sensación de tres de los cinco sentidos: audición, visión y tacto (este último brevemente en el Cap. 2). En el Cap. 1 se realiza alguna conexión entre la física de la sensación, la bioquímica de la sensación y la percepción (psicofísica). Los sentidos del gusto y del olfato son puramente químicos, con poca base en la física (aparte de que la química de las interacciones moleculares en cada uno de ellos son aplicaciones claras de la física), y no se cubren, excepto por una breve discusión de las propiedades eléctricas de las neuronas sensoriales del gusto y del olfato en el Cap. 12. El énfasis en todo momento está en cómo la física puede explicar el funcionamiento del cuerpo en circunstancias normales e inusuales. Nos ocuparemos del cuerpo humano con sus coberturas corporales comunes: calzado para minimizar el estrés durante el movimiento (Cap. 4), ropa para regular la pérdida de calor (Cap. 6) y lentes correctivos para mejorar la visión (Cap. 11).

Los capítulos están diseñados más para abordar áreas específicas de la física que partes o sistemas específicos del cuerpo. Es difícil construir capítulos con divisiones nítidas porque se necesitan diferentes áreas de la física para comprender muchos componentes del cuerpo. Por ejemplo, para comprender la física del corazón, es necesario abordar su papel en la circulación (Cap. 8), la acción de los músculos (Cap. 5, que se centra más en el músculo esquelético que en el músculo cardíaco, que es bastante similar) y las señales eléctricas generadas por el corazón (Cap. 12).

Este texto concluye con cinco apéndices. El Apéndice A ofrece una visión general de los símbolos y las unidades, y hace referencia a las tablas de unidades presentadas en los capítulos. El Apéndice B enumera las figuras y tablas que describen las características principales de la información anatómica y antropométrica humana, que se utilizan a lo largo de este texto. Los tipos de ecuaciones diferenciales utilizadas en el texto se revisan en el Apéndice C. Estas mismas ecuaciones diferenciales se utilizan en todo el texto en modelos mecánicos, de flujo de fluidos y eléctricos; las conexiones entre estos modelos se realizan en el Apéndice D. El Apéndice E intenta definir el campo de la biofísica y conecta los contenidos de este texto con este campo.

Este texto se ha desarrollado a partir de las notas de clase del autor preparadas para el curso \textit{Física del Cuerpo Humano}, que es una materia optativa restringida de «nivel profesional» que él diseñó, cursada principalmente por estudiantes de primer y segundo año de pregrado en la Fu Foundation School of Engineering and Applied Science de la Universidad de Columbia. Este curso fue diseñado de modo que pudiera ser tomado por todos los estudiantes de primer año en su segundo semestre (en conjunto con la física y el cálculo del segundo semestre).

El autor suele cubrir los Caps. 1 al 10 con cierto detalle y los Caps. 11 al 13 con menos detalle en un semestre completo.

Se pueden dictar cursos a diferentes niveles, incluidos cursos de pregrado de nivel medio y superior, incluyendo o excluyendo deliberadamente temas más detallados y avanzados en el texto y los problemas. Dependiendo del nivel de profundidad deseado, el material de aproximadamente la mitad o de todos los capítulos puede cubrirse en un semestre.

Este texto también puede utilizarse como volumen complementario en cursos introductorios de física, y ayudar a los estudiantes de pregrado de premédica a aprender y repasar física. También puede servir como texto en cursos introductorios de ingeniería biomédica o de materiales médicos. Los estudiantes de medicina interesados en un enfoque más cuantitativo de la fisiología y aquellos que realizan investigación médica también pueden apreciar los enfoques adoptados aquí.

Al final de cada capítulo se presentan muchos problemas que van desde problemas sencillos hasta otros más avanzados (estos últimos se denominan como tales). Varios problemas tienen múltiples partes, y solo se pueden asignar algunas de estas partes. Las respuestas a los problemas seleccionados se proporcionan después de los apéndices.

Por lo general, se utilizan unidades del SI (MKS, m-kg-s); cuando resulta más conveniente, se utilizan otras unidades métricas, incluidas las unidades CGS (cm-g-s) y unidades métricas mixtas. A veces se utilizan deliberadamente unidades del sistema inglés FPS (ft-lb-s) para hacer una conexión con el mundo real (al menos en países como EE. UU. y el Reino Unido). Por ejemplo, sería extraño escuchar a un locutor de béisbol decir: «Este lanzador realmente está tirando fuego. La pistola de radar cronometró su último lanzamiento a $43,8\text{ m/s}$ (o $158\text{ km/h}$), en lugar de $98\text{ mph}$». Sería aún más extraño escuchar a un locutor de fútbol (es decir, fútbol americano) decir: «Tienen primero y $9,144$ por avanzar», en referencia a $9,144\text{ m}$ en lugar de $10\text{ yd}$. Del mismo modo, sería extraño discutir la física del cuerpo en estos deportes, como al lanzar una pelota de béisbol, en unidades distintas de las habituales. Los ángulos se expresan en radianes, excepto cuando el uso de grados proporciona una imagen física más clara.

Varios textos excelentes cubren material que se superpone con los temas tratados aquí, cada uno con un enfoque diferente. Son recursos magníficos por derecho propio. \textit{Physics of the Body} de Cameron, Skofronick y Grant \cite{1st-6} abarca la mayoría de los temas de este libro y proporciona una excelente visión física. Está a un nivel de física inferior al utilizado aquí y deduce y presenta menos de las ecuaciones necesarias para un tratamiento más riguroso, pero proporciona una base básica muy buena en fisiología humana para no expertos. En cierto modo, el énfasis de \textit{The Human Machine} de Alexander \cite{1st-2} coincide con el nuestro, pero, de nuevo, las explicaciones son más cualitativas. El modo de pensamiento físico que presenta es impresionante. \textit{Physics with Examples from Medicine and Biology} de Benedek y Villars \cite{1st-3} es una serie de tres textos completos de física introductoria en los que se han elegido excelentes ejemplos y problemas relativos a la física del cuerpo. Los autores han tomado varios temas bastante complejos y los han hecho completamente comprensibles. Muchos otros textos de física general de primer año de uso común hoy en día contienen varios ejemplos y problemas de capítulo relacionados con el cuerpo. \textit{Intermediate Physics for Medicine and Biology} de Hobbie \cite{1st-14} es un texto más avanzado que enfatiza tanto la física como la fisicoquímica. \textit{Medical Physics and Biomedical Engineering} de Brown et al. \cite{1st-4} es un poco más avanzado y se centra también en áreas clásicas de la física médica, como la radiactividad y la instrumentación. Muchos de los problemas ilustrativos sobre biología humana y temas relacionados se han recopilado en los hermosos libros: \textit{Biomedical Applications of Introductory Physics} de Tuszynski y Dixon \cite{1st-28}, \textit{Physics in Biology and Medicine} de Davidovits \cite{1st-9}, \textit{Biophysics Problems: A Textbook with Answers} de Maróti, Berkes y Tölgyesi \cite{1st-17}, \textit{Physics for the Biological Sciences: A Topical Approach to Biophysical Concepts} de Hallett, Stinson y Speight \cite{1st-12}, y \textit{Topics in Classical Biophysics} de Metcalf \cite{1st-18}. Muchos de los temas de fisiología del ejercicio, como el metabolismo durante las actividades deportivas, se describen en términos elementales en \textit{Fox's Physiological Basis for Exercise and Sport} de Foss y Keteyian \cite{1st-10} y \textit{Physiology of Sport and Exercise} de Wilmore y Costill \cite{1st-30}. \textit{Basic Biomechanics of the Musculoskeletal System}, editado por Nordin y Frankel \cite{1st-20}, es una descripción general, completa y clara de la biomecánica de estructuras, articulaciones y movimiento. Las aplicaciones de la física a un nivel más molecular y celular, más en el dominio clásico de la biofísica, se describen en \textit{Biophysics: An Introduction} de Cotterill \cite{1st-8} y \textit{Biological Physics: Energy, Information} de Nelson \cite{1st-19}. La aplicación más general de la física a los animales se aborda en el apasionante y muy completo libro \textit{Zoological Physics: Quantitative Models, Body Design, Actions and Physical Limitations in Animals} de Ahlborn \cite{1st-1}. Todos estos textos son muy recomendados para obtener más detalles. Ellos, junto con los textos de anatomía y fisiología citados anteriormente, han contribuido a la preparación de este texto.

El autor agradece a las muchas personas que han realizado valiosos comentarios para contribuir a este libro, incluyendo a Marlene Arbo, Gerard Ateshian, Sarbajit Benerjee, Alex Breskin, Bill Burdick, Yi-Ting Chiang, Kevin Costa, Ted Ducas, Yossi Goffer, Daniel Herman, Jonathan Herman, Steven Heymsfield, Jeffrey Holmes, Mark Langill, Barclay Morrison III, Elizabeth Olson, Thomas Pedersen, Harry Radousky, Paul Sajda, Michael Sheetz y Samuel Sia. También le gustaría agradecer al sistema de bibliotecas de la Universidad de Columbia.

Este autor comenzó a escribir este texto cuando era becario Lady Davis durante su año sabático en la Universidad Hebrea de Jerusalén como invitado de Uri Banin, y agradece sinceramente este apoyo.

\vspace{1.5em}
\noindent Nueva York, NY \hfill Irving P. Herman \\
Noviembre de 2006

\vspace{2em}
\begin{thebibliography}{99}
\bibitem{1st-1} B.K. Ahlborn, \textit{Zoological Physics: Quantitative Models, Body Design, Actions and Physical Limitations in Animals} (Springer, Berlín, Heidelberg, Nueva York, 2004)
\bibitem{1st-2} R. McNeill Alexander, \textit{The Human Machine} (Columbia University Press, Nueva York, 1992)
\bibitem{1st-3} G.B. Benedek, F.M.H. Villars, \textit{Physics with Illustrative Examples from Medicine and Biology}, tres volúmenes, 2.ª ed. (AIP/Springer, Nueva York, Berlín, Heidelberg, 2000)
\bibitem{1st-4} B.H. Brown, R.H. Smallwood, D.C. Barber, P.V. Lawford, D.R. Hose, \textit{Medical Physics and Biomedical Engineering} (Institute of Physics, Filadelfia, 1999)
\bibitem{1st-5} R.F. Burton, \textit{Physiology by Numbers: An Encouragement to Quantitative Thinking} (Cambridge University Press, Cambridge, 1994)
\bibitem{1st-6} J.R. Cameron, J.G. Skofronick, R. Grant, \textit{Physics of the Body}, 2.ª ed. (Medical Physics, Madison, WI, 1999)
\bibitem{1st-7} E. Carson, C. Cobelli, \textit{Modelling Methodology for Physiology and Medicine} (Academic, San Diego, 2001)
\bibitem{1st-8} R.M.J. Cotterill, \textit{Biophysics: An Introduction} (Wiley, Nueva York, 2002)
\bibitem{1st-9} P. Davidovits, \textit{Physics in Biology and Medicine}, 2.ª ed. (Elsevier/Academic, San Diego, 2001)
\bibitem{1st-10} M.L. Foss, S.J. Keteyian, \textit{Fox's Physiological Basis for Exercise and Sport}, 6.ª ed. (McGraw-Hill, Boston, 1998)
\bibitem{1st-11} A.C. Guyton, J.E. Hall, \textit{Human Physiology and Mechanisms of Disease}, 6.ª ed. (Saunders, Filadelfia, 1997)
\bibitem{1st-12} F.R. Hallett, R.H. Stinson, P.A. Speight, \textit{Physics for the Biological Sciences: A Topical Approach to Biophysical Concepts} (Methuen/Chapman \& Hall, Toronto, 1982)
\bibitem{1st-13} J.L. Hargrove, \textit{Dynamic Modeling in the Health Sciences} (Springer, Berlín, Heidelberg, Nueva York, 1998)
\bibitem{1st-14} R.K. Hobbie, \textit{Intermediate Physics for Medicine and Biology}, 3.ª ed. (AIP, Nueva York, 1997)
\bibitem{1st-15} F.C. Hoppensteadt, C.S. Peskin, \textit{Modeling and Simulation in Medicine and the Life Sciences}, 2.ª ed. (Springer, Berlín, Heidelberg, Nueva York, 2002) (La primera edición de 1992 se llamó \textit{Mathematics in Medicine and the Life Sciences})
\bibitem{1st-16} E.N. Marieb, \textit{Human Anatomy \& Physiology}, 5.ª ed. (Benjamin Cummings, San Francisco, 2001)
\bibitem{1st-17} P. Maróti, L. Berkes, F. Tölgyesi, \textit{Biophysics Problems: A Textbook with Answers} (Akadémiai, Budapest, 1998)
\bibitem{1st-18} H.J. Metcalf, \textit{Topics in Classical Biophysics} (Prentice-Hall, Englewood Cliffs, NJ, 1980)
\bibitem{1st-19} P. Nelson, \textit{Biological Physics: Energy, Information, Life} (Freeman, Nueva York, 2003)
\bibitem{1st-20} M. Nordin, V.H. Frankel (eds.), \textit{Basic Biomechanics of the Musculoskeletal System}, 3.ª ed. (Lippincott Williams \& Wilkins, Filadelfia, 2001)
\bibitem{1st-21} P.E. Pack, \textit{Anatomy and Physiology} (Cliff Notes, Lincoln, NE, 1997)
\bibitem{1st-22} H.D. Patton, A.F. Fuchs, B. Hille, A.M. Scher, R. Steiner (eds.), \textit{Textbook of Physiology}, 21.ª ed. (Saunders, Filadelfia, 1989)
\bibitem{1st-23} R.R. Seeley, T.D. Stephens, P. Tate, \textit{Anatomy \& Physiology}, 6.ª ed. (McGraw-Hill, Boston, 2003)
\bibitem{1st-24} L. Sherwood, \textit{Human Physiology: From Cells to Systems}, 5.ª ed. (Thomson Brooks/Cole, Belmont, CA, 2004)
\bibitem{1st-25} D. Shier, J. Butler, R. Lewis, \textit{Hole's Human Anatomy and Physiology}, 9.ª ed. (McGraw-Hill, Boston, 2002)
\bibitem{1st-26} A. Stalheim-Smith, G.K. Fitch, \textit{Understanding Human Anatomy and Physiology} (West, St. Paul, MN, 1993)
\bibitem{1st-27} G.J. Tortora, S.R. Grabowski, \textit{Principles of Anatomy and Physiology}, 10.ª ed. (Wiley, Nueva York, 2003)
\bibitem{1st-28} J.A. Tuszynski, J.M. Dixon, \textit{Biomedical Applications of Introductory Physics} (Wiley, Nueva York, 2002)
\bibitem{1st-29} J.B. West (ed.), \textit{Best and Taylor's Physiological Basis of Medical Practice}, 12.ª ed. (Williams \& Wilkins, Baltimore, 1991)
\bibitem{1st-30} J.H. Wilmore, D.L. Costill, \textit{Physiology of Sport and Exercise}, 3.ª ed. (Human Kinetics, Champaign, IL, 2004)
\end{thebibliography}

% Control de Calidad: 20 párrafos origen -> 20 párrafos salida. Mapeo 1:1 verificado. Sin omisiones.

\end{document}

